% GENERATED FILE. Edit canonical lessons, then run npm run generate:books.
% canonical-source-hash: fnv1a64:55843a421b81aaae
% canonical-lessons: RU-C106-dessert, RU-C106-yoghurt, RU-C106-ice-cream, RU-C106-sausage, RU-C106-peach

\chapter{Dessert, Yoghurt, Ice cream, Sausage, Peach}
\label{ch:ru-dessert-yoghurt-ice-cream-sausage-peach}
\hlchaptermodality{106}

\begin{chapteropening}
\textbf{By the end of this chapter:} \emph{I can say a dessert, a yoghurt, ice cream, sausage and a peach.}

Complete the last lesson of chapter 106: I can say a dessert, a yoghurt, ice cream, sausage and a peach.
\end{chapteropening}

\section[\texorpdfstring{\ru{десерт} (desért)}{десерт (desért)}]{\ru{десерт} (desért) — a dessert}

\label{lesson:RU-C106-dessert}



\begin{quote}
Before the new one: say the Russian for a sandwich, then the Russian for a sauce.
\end{quote}

\subsection*{You'll want to know: \ru{десерт}}
\textbf{\ru{десерт}} (\emph{desért}) — \textquotedblleft{}a dessert\textquotedblright{}.

\begin{cousinweb}
\textbf{\ru{десерт}} came through German \emph{Dessert} or straight from French \emph{dessert} (Vasmer), from \emph{desservir}, \textquotedblleft{}to clear the table\textquotedblright{}.
\end{cousinweb}

\begin{grammarlens}[title={What you've built}]
One more word for this chapter.
\end{grammarlens}

\subsection*{Your turn}
\begin{itemize}
\raggedright
  \item \emph{Say it:} \emph{desért}
  \item \emph{Say it:} \emph{desért}, once more
  \item \emph{Say it:} say \emph{desért}
  \item \emph{Recall:} say the Russian for a sandwich, then the Russian for a sauce, then say \emph{desért} again
\end{itemize}

\begin{tcolorbox}[breakable,colback=teal!4,colframe=teal!35!black,title={Before you move on}]
What does \ru{десерт} mean? (\textquotedblleft{}A dessert\textquotedblright{}.) Say it once more.
\end{tcolorbox}
\section[\texorpdfstring{\ru{йогурт} (yógurt)}{йогурт (yógurt)}]{\ru{йогурт} (yógurt) — a yoghurt}

\label{lesson:RU-C106-yoghurt}



\begin{quote}
Before the new one: say the Russian for a sauce, then the Russian for a dessert.
\end{quote}

\subsection*{You'll want to know: \ru{йогурт}}
\textbf{\ru{йогурт}} (\emph{yógurt}) — \textquotedblleft{}a yoghurt\textquotedblright{}.

\begin{cousinweb}
\textbf{\ru{йогурт}} comes from Turkish \emph{yoğurt}, \textquotedblleft{}yoghurt\textquotedblright{}, through the western European languages.
\end{cousinweb}

\begin{grammarlens}[title={What you've built}]
One more word for this chapter.
\end{grammarlens}

\subsection*{Your turn}
\begin{itemize}
\raggedright
  \item \emph{Say it:} \emph{yógurt}
  \item \emph{Say it:} \emph{yógurt}, once more
  \item \emph{Say it:} say \emph{yógurt}
  \item \emph{Recall:} say the Russian for a sauce, then the Russian for a dessert, then say \emph{yógurt} again
\end{itemize}

\begin{tcolorbox}[breakable,colback=teal!4,colframe=teal!35!black,title={Before you move on}]
What does \ru{йогурт} mean? (\textquotedblleft{}A yoghurt\textquotedblright{}.) Say it once more.
\end{tcolorbox}
\section[\texorpdfstring{\ru{мороженое} (morózhenoye)}{мороженое (morózhenoye)}]{\ru{мороженое} (morózhenoye) — ice cream}

\label{lesson:RU-C106-ice-cream}



\begin{quote}
Before the new one: say the Russian for a dessert, then the Russian for a yoghurt.
\end{quote}

\subsection*{You'll want to know: \ru{мороженое}}
\textbf{\ru{мороженое}} (\emph{morózhenoye}) — \textquotedblleft{}ice cream\textquotedblright{}.

\begin{cousinweb}
\textbf{\ru{мороженое}} is an adjective standing alone: \textquotedblleft{}frozen (food)\textquotedblright{}, from \emph{\ru{морозить}}, \textquotedblleft{}to freeze\textquotedblright{}, and \emph{\ru{мороз}}, \textquotedblleft{}frost\textquotedblright{}, which Vasmer relates to Albanian \emph{mardhem}, \textquotedblleft{}I freeze\textquotedblright{}.
\end{cousinweb}

\begin{grammarlens}[title={What you've built}]
One more word for this chapter.
\end{grammarlens}

\subsection*{Your turn}
\begin{itemize}
\raggedright
  \item \emph{Say it:} \emph{morózhenoye}
  \item \emph{Say it:} \emph{morózhenoye}, once more
  \item \emph{Say it:} say \emph{morózhenoye}
  \item \emph{Recall:} say the Russian for a dessert, then the Russian for a yoghurt, then say \emph{morózhenoye} again
\end{itemize}

\begin{tcolorbox}[breakable,colback=teal!4,colframe=teal!35!black,title={Before you move on}]
What does \ru{мороженое} mean? (\textquotedblleft{}Ice cream\textquotedblright{}.) Say it once more.
\end{tcolorbox}
\section[\texorpdfstring{\ru{колбаса} (kolbasá)}{колбаса (kolbasá)}]{\ru{колбаса} (kolbasá) — sausage}

\label{lesson:RU-C106-sausage}



\begin{quote}
Before the new one: say the Russian for a yoghurt, then the Russian for ice cream.
\end{quote}

\subsection*{You'll want to know: \ru{колбаса}}
\textbf{\ru{колбаса}} (\emph{kolbasá}) — \textquotedblleft{}sausage\textquotedblright{}.

\begin{cousinweb}
\textbf{\ru{колбаса}} — recorded from 1280; origin disputed. Some relate it to \emph{\ru{колоб}}, \textquotedblleft{}a round loaf\textquotedblright{}; Vasmer finds grounds to suspect a loan and names Turkish \emph{külbastı}, \textquotedblleft{}grilled meat\textquotedblright{}, as a possible source.
\end{cousinweb}

\begin{grammarlens}[title={What you've built}]
One more word for this chapter.
\end{grammarlens}

\subsection*{Your turn}
\begin{itemize}
\raggedright
  \item \emph{Say it:} \emph{kolbasá}
  \item \emph{Say it:} \emph{kolbasá}, once more
  \item \emph{Say it:} say \emph{kolbasá}
  \item \emph{Recall:} say the Russian for a yoghurt, then the Russian for ice cream, then say \emph{kolbasá} again
\end{itemize}

\begin{tcolorbox}[breakable,colback=teal!4,colframe=teal!35!black,title={Before you move on}]
What does \ru{колбаса} mean? (\textquotedblleft{}Sausage\textquotedblright{}.) Say it once more.
\end{tcolorbox}
\section[\texorpdfstring{\ru{персик} (pérsik)}{персик (pérsik)}]{\ru{персик} (pérsik) — a peach}

\label{lesson:RU-C106-peach}



\begin{quote}
Before the new one: say the Russian for ice cream, then the Russian for sausage.
\end{quote}

\subsection*{You'll want to know: \ru{персик}}
\textbf{\ru{персик}} (\emph{pérsik}) — \textquotedblleft{}a peach\textquotedblright{}.

\begin{cousinweb}
\textbf{\ru{персик}} probably came through Middle Low German \emph{persik} from Latin \emph{(mālum) persicum}, \textquotedblleft{}the Persian apple\textquotedblright{} (Vasmer).
\end{cousinweb}

\begin{grammarlens}[title={What you've built}]
One more word for this chapter.
\end{grammarlens}

\subsection*{Your turn}
\begin{itemize}
\raggedright
  \item \emph{Say it:} \emph{pérsik}
  \item \emph{Say it:} \emph{pérsik}, once more
  \item \emph{Say it:} say \emph{pérsik}
  \item \emph{Recall:} say the Russian for ice cream, then the Russian for sausage, then say \emph{pérsik} again
\end{itemize}

\begin{tcolorbox}[breakable,colback=teal!4,colframe=teal!35!black,title={Before you move on}]
What does \ru{персик} mean? (\textquotedblleft{}A peach\textquotedblright{}.) Say it once more.
\end{tcolorbox}
